\section{The unramified lifting exponent on a split middle-range declaration}

% BEGIN L81S1
\paragraph{L81S1. The place tests, coefficient extension and height.}
Use F78D2's hub $X$. Fix a twisted $i$ with $p=\ell_i$ outside the two
type primes, in $\mathrm{RangeAD}(L,X.\mathrm{exc},p)$ and split in the
host classification. At a column the matrix family is
$R_x\chi_1^{-1}$, for every $x\in X_i$; at a row it is
$(R_x\chi_1^{-1})^\vee(1)$. Both use the same stored integer
\[
 h=\lfloor h_B\rfloor,\quad h_B=b+S_c+\eta+t+2,\qquad
 e_S=X.e_S(i)=\left\lceil23\left(2+2h/p+
                         \log_p(2h+1)\right)\right\rceil .
\]
Indeed the canonical row data set $e_S=e_A(L,23,p)$ at split
primes. The general column test takes a function of the place
$w$. The split declaration uses the constant function $e_S$.

Fix any $w\mid p$ and any tested good coefficient extension $L'$
with integers $O$. M81B1 gives the unramified field $K=K_0$,
the good level-one layer and odd $p$. Away from the type primes,
the hub facts give crystallinity of every $R_x$ as in S80J3.
For the first character one may use the rank-one construction
even if there is no good point at $i$. The block dictionary makes
the infinitesimal character of the first Hecke character integral
after the specified shift (its entry is the negative first jump).
The rank-one case of the given reciprocity field produces its
continuous Galois character with the prescribed Weil--Deligne
parameters. Its Frobenius values agree with those of $\chi_1$
outside the finitely many excluded places, so Chebotarev and
continuity identify the two characters. At $w\mid p$ the character
is de Rham, and its Weil--Deligne parameter is unramified by the
hub's $\mathrm{MuUnr}$ condition and unramified character twist.
The bridge's de Rham/unramified Weil--Deligne criterion gives
crystallinity. Extending coefficients, dualizing and tensoring
crystalline bases gives both crystalline matrix families over $L'$.

S80J1--S80J2 give $|z|\le h_J=b+S_c+\eta+t$ for every jump
of the twisted family; the dual twists have jumps $-1-z$.
Both lists therefore lie in $[-h,h]$. A test in $\mathrm{LiftAt}$
uses an onto lattice map $f:T_1\to T_2$ with rational representations
subquotients of finite sums from the indicated family. Crystalline
comparison and strictness on subquotients (H80M1) give
crystallinity and the same jump window for the actual lattices
and their kernel. The following proof is uniform in their ranks.
% END L81S1

% BEGIN L81S2
\paragraph{L81S2. Tame invariants and the two integral Wach lifts.}
Put $s=p-1$ and $\Delta=\mu_{p-1}(\mathbb Z_p)$, acting by
$[d](X)=(1+X)^d-1$. Apply A81R3's integral coordinate
$Z=s^{-1}\sum_{d\in\Delta}d^{-1}[d](X)$, which transforms by
the character $d$ and has linear coefficient one. Set initially
$v_0=Z^s$. At any primitive $p$th root $\zeta$ the value
$v_0(\zeta-1)$ is independent of the root by tame transitivity.
It belongs to $\mathbb Q_p$ (the coefficients are in $\mathbb Z_p$)
and has valuation one, since $Z/X$ is a unit. Rescale $v_0$ by
the $\mathbb Z_p$-unit $-p/v_0(\zeta-1)$, calling the result $v$.
Then $E=v+p$ is a unit times $Q=\Phi_p(1+X)$: divide by the
distinct factors $X-(\zeta-1)$ and compare the residue order $s$.
The invariant series on each letter are $O[[v]]$ and
$\varphi(v)=vE^s$ times a unit. The procyclic generator $\gamma$,
$u=\chi(\gamma)$, $v_p(u-1)=1$, satisfies $\gamma(v)-v\in vE O[[v]]$.
Indeed the difference divided by $v$ vanishes at every primitive
root, so is divisible by $Q$, and the quotient by $E$ is invariant.
Consequently $\gamma(E)/E$ is a unit congruent to one modulo $v$.

Write $N(T)$ here for the product over letters of the
$\Delta$-invariants of the given Wach module. W79H1--W79H3
at $X=0$ show that the tame group acts trivially modulo $X$
on the original module: the regular Taylor action is substitution
on coefficients with fixed crystalline constants. Averaging a
Wach basis as in A81R7 then gives an invariant basis over
$O[[v]]$, with $\gamma$ acting trivially modulo $v$.
The height bound $E^h\varphi N\subset N$ follows from W79H4.

Set $w(k)=k+v_p(k!)$, $A=\lceil2h/s\rceil$ and
$b_W=w(2A-1)$. Then $p^{b_W}N(T_2)\subseteq fN(T_1)$.
Indeed A81R8's Frobenius coefficient argument uses precisely
$\varphi(X)=XQ$, the two Wach heights $h$, and surjectivity on
the residual \'etale modules. All three hold here, so $X^{2h}$
kills the full Wach cokernel. Thus $v^A$ kills its invariant
cokernel by averaging a preimage. On successive $v^j$ quotients
the generator acts by $u^{sj}$. The two polynomials with root
indices $0\le j<A$ and $A\le j<2A$ kill $v^A N(T_2)$
modulo $v^A fN(T_1)$ as in A81R9. Lagrange interpolation on
these $2A$ roots costs $p^{b_W}$: each denominator has valuation
$2A-1+v_p(k!)+v_p((2A-1-k)!)\le b_W$ at index $k$.
Cancellation of $v^A$ gives the assertion. The same estimate
holds on integral extensions by a trivial quotient, whose new
jump is zero. Here $h\ge1$ by the family size conditions.

A bounded target pair has $y\in E^{-h}N(T_2)$, $z\in vN(T_2)$
and $(\gamma-1)y=(\varphi-1)z$. For every integral finite
target cocycle $c$ there exists such a pair representing
$p^{b_W}[c]$ on restriction to $F=K(\mu_p)$: lift $p^{b_W}$
times the fixed generator in the invariant Wach lattice of the
crystalline extension defined by $c$, then take its Frobenius
and generator differences. The kernel identifications of A81R10
and B81L1 apply here by their height-one argument (strict
primitive Taylor lattice, full rank at zero and on residual
\'etale modules); taking $\Delta$-invariants preserves kernels.
Thus these differences have the stated bounds. On the Witt tensor
the lifted generator minus $p^{b_W}$ times the chosen quotient
lift gives the two exact primitive identities of B81L3.
The rational finite-class surjectivity proof there, over $K$,
uses strict filtered exactness and the crystalline dimension
formula. Its scaled integral extension lattice, followed by the
invariant Wach lift just proved, therefore gives a bounded
source pair above $p^k(y,z)$ for some $k$. This auxiliary power
does not enter the final estimate.
% END L81S2

% BEGIN L81S3
\paragraph{L81S3. The divided difference on the kernel pole lattice.}
For the kernel write $N=N(T_0)$, $N_Q=E^{-h}N$ and
$\delta=(\gamma-1)/v$. Both $N,N_Q$ are stable under $\delta$:
the congruences modulo $v$ on $N$ and $\gamma(E)/E$ prove
this directly. Write $U\subset[-h,h]$ for the source's jump set
and $J=((U+\{1,\ldots,h\})\setminus\{0\})\cup\{1,\ldots,h\}$.
Set $c_j=(1-u^{-j})/p\in\mathbb Z_p$ and
$P(B)=\prod_{j\in J}(B-c_j)$. Then
\[
                    \delta P(\delta)N_Q\subseteq N .
\]
The quotient $N_Q/N$ is finite free over $O$ by division by
the monic polynomial $(v+p)^h$, so it is enough to test its
rational Taylor jet at $v=-p$ on every letter. The Taylor
coordinate $1+X=\zeta\exp(Y)$ is equivalent there, since
$E\sim Q$ has a simple zero at $\zeta-1$. The invariant basis
from L81S2 is also a basis before taking invariants, so the
filtered Taylor lattice in these coordinates is exactly that of
W79H2--W79H3. In a fixed filtration-adapted period basis each
line of jump $a$ has the basis $Y^{-a}$ times its period, and the
$h$ pole degrees have exponents $-a-1,\ldots,-a-h$.
On the degree $-j$ term, substitution by $uY$ and multiplication
by $v(Y)^{-1}$ in $\delta$ give diagonal coefficient
$(u^{-j}-1)/(-p)=c_j$ and only higher-degree additional terms.
This triangular operator preserves each period line, where its
diagonal entries are pairwise distinct over $L'$ because $u$
has infinite order. It is diagonalizable on that line.
Strictness puts its $a$ among the jumps of $T_1$, so all these
diagonal entries are roots of $BP(B)$, including the possible
zero. This proves the inclusion without multiplying the constant
denominators from the Taylor change of basis.

Put $k_0=\lfloor h/s\rfloor$, $C=\prod_{j=1}^{k_0}(u^{sj}-1)$.
The graded $v^j$ action gives
$CvN\subseteq(\gamma-1)vN+v^{k_0+1}N$ by the polynomial
argument of B81L4. On $v^{k_0+1}N$ the operator $1-\varphi$
is invertible by geometric series: applying Frobenius contributes
an extra $E^{s(k_0+1)-h}$ with positive exponent, and $(p,v)$
is Frobenius-stable. Both calculations include $k_0=0$.
% END L81S3

% BEGIN L81S4
\paragraph{L81S4. Correcting the divided integral defect.}
On $N_Q\oplus vN$ form the divided defect
\[
           D(y,z)=\delta y-(\varphi-1)z/v\in N_Q .
\]
The inclusion follows from $\varphi(v)/v\sim E^s$. Suppose
$p^k\omega=D(y,z)$ with $\omega\in N_Q$. Write
$Cz=(\gamma-1)x+z_1$, $x\in vN$, $z_1\in v^{k_0+1}N$
by L81S3. Commuting before division by $v$ gives
\[
 Cp^k\omega=\delta\big(Cy-(\varphi-1)x\big)
                      -(\varphi-1)z_1/v .
\]
The first argument lies in $N_Q$ and the last term in $N$.
L81S3 and $\delta N\subset N$ give $Cp^k P(\delta)\omega\in N$.
Torsion freeness of $N_Q/N$ yields $P(\delta)\omega\in N$.
For each $n\in N$, apply the regulator to $vn$ to write
$Cvn=(\gamma-1)x'+z'_1$ with $x'\in vN$ and
$z'_1\in v^{k_0+1}N$. Invert $1-\varphi$ on the last term
and divide the equation by $v$. Thus $Cn\in\mathrm{im}\,D$.
It follows that $CP(\delta)\omega$ lies in this image.
Since $P(B)-P(0)$ is divisible by $B$, its contribution lies in
$\delta N_Q$ and hence in $\mathrm{im}\,D$. Therefore
$p^{d_U}\omega\in\mathrm{im}\,D$, with
\[
             d_U=v_p(CP(0))=w(k_0)+\sum_{j\in J}v_p(j).
\]
Here $v_p(u^{sj}-1)=1+v_p(j)$, whereas division by $p$ in
$c_j$ leaves only $v_p(j)$.

For L81S2's target pair lift $p^{b_W}E^hy$ and $p^{b_W}z/v$
through the invariant Wach map, then divide or multiply back.
The divided defect of these lifts lies in the kernel $N_Q$;
the auxiliary bounded lift over $p^k(y,z)$ makes its $p^k$
multiple a kernel defect by subtraction as in B81L5.
The calculation above corrects it at exactly the allowed cost
$p^{d_U}$. We obtain an invariant bounded source pair mapping
to $p^{b_W+d_U}(y,z)$. Its restriction to $G_F$ has poles at
most $Q^h$ and second coordinate in $XN$, so B81L2 constructs
a crystalline integral extension there. B81L3's Witt identity
shows its class maps to $p^{2b_W+d_U}[c]|_{G_F}$.
Corestriction and division by the unit $[F:K]=p-1$ yield an
integral finite lift of $p^{2b_W+d_U}[c]$ over $K$.

Finally $J\subset[-h,2h]\setminus\{0\}$ gives
$\sum_{j\in J} v_p(j)\le v_p(h!)+v_p((2h)!)\le3h/s$.
Legendre's bound gives $w(m)\le mp/s$, $2A-1\le4h/s+1$,
and $k_0\le h/s$. Consequently
\[
 2b_W+d_U\le(9p/s^2+3/s)h+2p/s
            \le25h/p+3\le e_S ,
\]
using $p/s\le3/2$ and L81S1's ceiling at the last step.
Multiplication by the remaining power proves both
$\mathrm{LiftExponentColumn}$ and $\mathrm{LiftExponentDualPlaces}$
at the stored $e_S$ on every guarded split index. This establishes
the two lifting fields of the split declaration; the inertia and
gluing conditions for its global index inequality are separate.
% END L81S4
